% Requires the ordinary theorem environments, amsmath/amssymb, and cleveref.
% External proof obligations: lem:clean-support-regularization and
% lem:five-input-obstruction. The latter must be proved in the same paper
% if the resulting proof is to be advertised as self-contained.

\subsection{A restriction with fixed adjacent pairs}

Write $(a)_r=a(a-1)\cdots(a-r+1)$. After normalization, the network has the
form
\begin{equation}\label{eq:fixed-pair-model}
 F(x)=\ell_0(x)+\sum_{j=1}^{h_1}d_j\sigma(\langle a_j,x\rangle)
             +\sum_{k=1}^{h_2}v_k\sigma(g_k(x)),
 \qquad
 g_k(x)=\ell_k(x)+\sum_{j=1}^{h_1}w_{kj}\sigma(\langle a_j,x\rangle),
\end{equation}
where all first-layer rows and the inner linear forms $\ell_k$, $k\ge1$,
are translation invariant. The output linear form $\ell_0$ carries the
unit translation slope of MAX. The auxiliary linear forms and shallow
output terms in this intermediate model do not change the architecture
in the theorem statement. In particular, the original second-layer index
set is retained.

Put $r_j=|\operatorname{supp}(a_j)|$ and
\[
 M=|\{j:r_j\ge3\}|,
 \qquad L=\sum_j(r_j)_4,
 \qquad Q=\sum_{j:r_j\ge3}r_j^4.
\]
Thus $M,L\le Q$. Every row of support two is a scalar multiple of
$e_u-e_v$. Using $\sigma(-t)=\sigma(t)-t$ to give all such rows a fixed
orientation, define
\[
 c_{k,uv}=\sum_{j:\operatorname{supp}(a_j)=\{u,v\}}
                  w_{kj}|a_{j,u}|,
 \qquad
 E(G_k)=\{\{u,v\}:c_{k,uv}\ne0\}.
\]
These coefficients are summed before the graph is defined, so signed
cancellation between parallel contributions is included.

Fix $n\ge10$, put $m=n-5$, and take an integer $1\le J\le m/2$.
Choose a uniform permutation of the $n$ coordinate labels. Its first five
labels form $S$; the others occupy ordered tail positions $1,\ldots,m$.
The selected coordinates will have value zero. Pair tail positions
\[
 (1,2),(3,4),\ldots,(2J-1,2J),
\]
and let $P$ be the resulting matching of coordinate labels. Position zero
is an anchor with value zero. The gap between positions $r-1$ and $r$ is
$T^{\kappa(r)}$, where $\kappa$ is a fixed bijection of $\{1,\ldots,m\}$
such that
\[
 \{\kappa(2),\kappa(4),\ldots,\kappa(2J)\}=\{1,\ldots,J\}.
\]
Thus all paired gaps have smaller exponents than every gap between blocks,
including the anchor gap. The value at tail position $t$ is
$y_t=-\sum_{r\le t}T^{\kappa(r)}$.

There is one sufficiently large $T$ that works for every sample and every
neuron. To see this, each first-row value is a polynomial in $T$, so its
sign eventually stabilizes. After all these finitely many signs stabilize,
each second preactivation is also a fixed polynomial. Choose $T$ beyond
all roots of the nonzero polynomials. For this common choice,
\begin{equation}\label{eq:zero-tail-gradient}
 g_k(y)=0\quad\Longrightarrow\quad
 \text{the coefficient vector of $g_k$ in the tail variables is zero}.
\end{equation}
Indeed, if that vector is $b$, the gap polynomial is
$-\sum_r(\sum_{t\ge r}b_t)T^{\kappa(r)}$. Distinct exponents force every
suffix sum to vanish, and hence $b=0$. A first row supported on $S$ is
identically zero on the tail subspace and causes no difficulty here.

Declare a first-row support of size at least three to be \emph{pair-bad}
if it contains either the endpoints of two pairs in $P$, or the endpoints
of one pair in $P$ and two labels of $S$. Let $B_{\rm pair}$ be the union
of these events. Let $B_{\rm selected}$ be the event that a first-row
support of size at least three is wholly contained in $S$.

\begin{lemma}\label{lem:fixed-pair-errors}
The two bad events satisfy
\[
 \mathbb P(B_{\rm selected})\le M\frac{(5)_3}{(n)_3},
 \qquad
 \mathbb P(B_{\rm pair})\le
 \frac{L}{(n)_4}\left[\binom J2+J\binom52\right].
\]
\end{lemma}
\begin{proof}
For the first bound, fix three labels in each eligible support. Their
probability of belonging to $S$ is $(5)_3/(n)_3$.
For the second, there are $\binom J2$ choices of two pair slots, and
$J\binom52$ choices of a pair slot and two selected positions. For each
choice, all four labels lie in a support of size $r$ with probability
$(r)_4/(n)_4$. Sum over choices and rows.
\end{proof}

\begin{lemma}\label{lem:fixed-pair-coins}
Fix $S$, the unordered labels in each pair, and all remaining tail labels.
If $X_k=|P\cap E(G_k)|$, then
\[
 \mathbb P\bigl(g_k(y)=0,\ B_{\rm pair}^{\mathsf c}
       \mid\text{these unordered data}\bigr)\le2^{-X_k}.
\]
\end{lemma}
\begin{proof}
The pair orientations are independent fair bits. The bad event is
unchanged by these bits, so assume that the fixed data are good.

Consider a zero-sum first-row covector $a$. Reversing a pair changes only
the coefficient of that pair's small gap; all large-gap coefficients are
unchanged. If some large-gap coefficient is nonzero, its highest nonzero
term determines the sign for our common large $T$. Thus the activation
gate is fixed under every pair reversal.

If every large-gap coefficient is zero, differences of successive block
boundary sums show that every singleton tail coefficient is zero, the
two coefficients in each pair add to zero, and the selected coefficients
add to zero. A row with tail support therefore contains an entire pair
with opposite nonzero coefficients. If it has support at least three,
it either contains a second such pair, or has a nonzero selected
coefficient. In the latter case it has at least two selected coefficients,
since their sum is zero. Both alternatives are pair-bad. Consequently,
on good data every non-braid row is either supported on $S$ or has a
fixed activation gate.

A braid row changes its gate only when its two labels form one of our
pairs. It follows that the tail coefficient vector of $g_k$ is
\[
 b_k+\sum_{i=1}^J\varepsilon_i c_{k,u_iv_i}(e_{u_i}-e_{v_i}),
 \qquad \varepsilon_i\in\{0,1\},
\]
with a fixed orientation convention. Its nonzero columns have disjoint
supports, hence are linearly independent. At most one assignment to the
$X_k$ relevant bits can make this vector zero. Apply
\Cref{eq:zero-tail-gradient}.
\end{proof}

% The following lemma is also supplied separately as matching/lemma.tex.
\begin{lemma}\label{lem:fixed-pair-matching}
Let $G$ have $e$ edges on $m\ge2$ vertices, and let $P$ be a uniformly
random matching of size $J$, with $0\le2J\le m$. Then
\[
 \mathbb E\,2^{-|P\cap E(G)|}\le\exp\!\left(-\frac{Je}{2m^2}\right).
\]
\end{lemma}
\begin{proof}
The cases $e=0$ or $J=0$ are immediate. Generate $P$ by a uniform
permutation, pairing its first $2J$ positions. For each graph edge $ab$,
swap $a$ into position one and then $b$ into position two. Every
permutation with these two fixed entries has exactly $m(m-1)$ preimages,
as one sees by undoing the swaps. After forgetting the pair order, the
result has the conditional matching law given $ab\in P$.

Put $X=|P\cap E(G)|$ and let $X_{ab}$ be the count after these swaps.
At most three pairs change, so $X_{ab}\le X+3$. With
$F(z)=\mathbb E z^X$ and $p=\mathbb P(ab\in P)=2J/[m(m-1)]$, we get
\[
 F'(z)=p\sum_{ab\in E(G)}\mathbb E z^{X_{ab}-1}
       \ge epz^2F(z),\qquad 0<z\le1.
\]
Integrating $F'/F\ge epz^2$ from $1/2$ to $1$ gives
$F(1/2)\le\exp(-7Je/[12m(m-1)])$, which implies the claim.
\end{proof}

\begin{proposition}\label{prop:fixed-pair-counting}
There is an absolute constant $D>0$ such that every normalized exact
$\operatorname{MAX}_n$ network as in \Cref{eq:fixed-pair-model} satisfies
\begin{equation}\label{eq:fixed-pair-counting}
 1-M\frac{(5)_3}{(n)_3}
  -\frac{L}{(n)_4}\left[\binom J2+J\binom52\right]
 \le \frac{D h_2}{J^2}
\end{equation}
for $n\ge10$ and $1\le J\le(n-5)/2$.
\end{proposition}
\begin{proof}
First fix a second neuron $k$ and put $e=|E(G_k)|$.
By a union bound over pairs of disjoint graph edges,
\[
 \mathbb P\bigl(G_k[S]\text{ contains two disjoint edges}\bigr)
 \le \frac{e^2}{2}\frac{(5)_4}{(n)_4}.
\]
Conditional on $S$, the tail pairs form a uniform $J$-matching on its
$m=n-5$ remaining labels. Deleting $S$ removes at most $5(n-1)$ edges.
The preceding two lemmas therefore bound the probability that this neuron
is zero, the data are pair-good, and $G_k[S]$ contains two disjoint edges by
\begin{equation}\label{eq:fixed-pair-neuron-cost}
 \frac{e^2}{2}\frac{(5)_4}{(n)_4}
 \exp\!\left(-\frac{J(e-5(n-1))}{2m^2}\right).
\end{equation}
Here the conditional bound is uniform in $S$; no independence of the
selected event and tail graph density is being assumed. The inequality
is also valid when $e-5(n-1)<0$.

Since $m\ge n/2$ and $J\le m/2$, the positive shift in the exponent is
bounded absolutely. Also $(n)_4\ge n^4/2$. On putting $u=Je/(2m^2)$,
\Cref{eq:fixed-pair-neuron-cost} is at most an absolute multiple of
$J^{-2}u^2e^{-u}$. The latter function of $u$ is bounded, so the neuron
cost is at most $D/J^2$ for an absolute $D$.

It remains to show that these neuron events cover every good sample.
Outside $B_{\rm selected}$, localization on the five tied coordinates
has only braid first-layer rows: genericity says a first-row covector
vanishes at the base point exactly when its support is contained in $S$.
 The localized network computes $\operatorname{MAX}_5$; its shallow output
 terms are a linear form and a sum of coordinate-difference ReLUs, as
 permitted by the structural lemma. By
\Cref{lem:five-input-obstruction}, one of its second preactivations has
two disjoint interaction edges. An original second neuron with nonzero
preactivation at the base point contributes only linear and shallow
first-layer terms, or zero, and supplies no retained outer ReLU.
Thus the witness is an original neuron with
$g_k(y)=0$, and its interaction graph on the selected coordinates is
exactly $G_k[S]$.

Every sample outside $B_{\rm selected}\cup B_{\rm pair}$ is therefore
covered by one of the $h_2$ events just bounded. Sum their probabilities
and apply \Cref{lem:fixed-pair-errors}.
\end{proof}

\begin{proof}[Proof of the existential-constant width bound]
Fix $C>0$ and replace it by $\max\{1,C\}$. By
\Cref{lem:clean-support-regularization}, the original biased network yields an exact
normalized model on $n=\lfloor N/2\rfloor$ inputs, with the same second
hidden width and $Q\le ACn^2$ for an absolute $A$.

Choose a sufficiently small absolute $\delta>0$, depending only on $A$,
and put $J=\lfloor\delta n/\sqrt C\rfloor$. For all sufficiently large
$n$ with $C$ fixed, $J\ge\delta n/(2\sqrt C)$ and $1\le J\le(n-5)/2$.
Since $M,L\le Q$, the two error terms in
\Cref{eq:fixed-pair-counting} are respectively
$O(C/n)$ and $O(CJ^2/n^2+CJ/n^2)$. The constants here are absolute.
Choose $\delta$ so that the quadratic contribution is at most $1/4$,
and then take $n$ large enough that the remaining terms have sum at most
$1/4$. The counting inequality gives
\[
 h_2\ge \frac{J^2}{2D}
      \ge \frac{\delta^2 n^2}{8DC}.
\]
As $n\ge N/3$ for $N\ge3$, this proves the desired quadratic bound
for all sufficiently large $N$, with $C$ fixed.

In the finite remaining range, $h_2\ge1$: an ordinary network with no
second-layer neurons has constant output and cannot compute
$\operatorname{MAX}_N$. Decrease the positive constant to cover these
finitely many values of $N$. This proves the existence of $c_C>0$ for
every fixed $C>0$.
\end{proof}
